\begin{mistake}{2026-07-20}{01}

\begin{problemcontent}
求双纽线 \(r^2 = a^2 \cos 2\theta \ (a > 0)\) 绕极轴旋转所成旋转曲面的面积。
\end{problemcontent}

\begin{solutioncontent}
由双纽线的对称性，取第一象限的曲线段（\(\theta \in [0, \frac{\pi}{4}]\)），其绕极轴（\(x\)轴）旋转生成的为右半部分曲面。总面积为右半部分的 \(2\) 倍。

对方程 \(r^2 = a^2 \cos 2\theta\) 两边对 \(\theta\) 求导得：
\[2r r' = -2a^2 \sin 2\theta \implies r' = -\frac{a^2 \sin 2\theta}{r}\]
计算弧长微元的平方：
\[r^2 + (r')^2 = r^2 + \frac{a^4 \sin^2 2\theta}{r^2} = \frac{r^4 + a^4 \sin^2 2\theta}{r^2}\]
将 \(r^4 = a^4 \cos^2 2\theta\) 代入分子：
\[r^2 + (r')^2 = \frac{a^4 \cos^2 2\theta + a^4 \sin^2 2\theta}{r^2} = \frac{a^4}{r^2}\]
故弧长微元 \(ds = \frac{a^2}{r} d\theta\)。

代入极坐标旋转曲面面积公式 \(S = \int 2\pi y ds = \int 2\pi (r \sin\theta) ds\)：
\[
\begin{aligned}
S &= 2 \times \int_{0}^{\pi/4} 2\pi (r \sin\theta) \frac{a^2}{r} d\theta \\
&= 4\pi a^2 \int_{0}^{\pi/4} \sin\theta d\theta \\
&= 4\pi a^2 \left[ -\cos\theta \right]_{0}^{\pi/4} \\
&= 4\pi a^2 \left( 1 - \frac{\sqrt{2}}{2} \right) = 2\pi a^2 (2 - \sqrt{2})
\end{aligned}
\]
\end{solutioncontent}

\mistakeknowledge{
\begin{itemize}
  \item \textbf{极坐标旋转面积公式}：绕极轴旋转 \(S = \int 2\pi (r\sin\theta) \sqrt{r^2 + (r')^2} d\theta\)，注意需保证 \(r\sin\theta \ge 0\)。
  \item \textbf{易错点}：切忌直接在 \([-\frac{\pi}{4}, \frac{\pi}{4}]\) 积分，第四象限 \(y<0\) 会导致面积抵消，必须利用对称性在第一象限积分并乘以倍数。
  \item \textbf{运算技巧}：遇 \(r^2 = f(\theta)\) 时，两边隐函数求导 \(2rr' = f'(\theta)\) 比直接开方求导更简便。
\end{itemize}}

\end{mistake}
